% year: תש"פ
% type: מבחן
% semester: א
% moed: א
% part: א
% number: 2
% points: 20
% categories: func-analysis, ivt, derivatives
% source: exams/מבחנים/תשפ/מועד א תשפ.pdf

\begin{question}
מצאו את הערך הגדול ביותר ואת הערך הקטן ביותר של הפונקציה $y=\sqrt[3]{x}(1+2x)$ בקטע $[-8,1]$.
\end{question}

\begin{hint}
המועמדים הם קצות הקטע, נקודות שבהן $y'=0$, ונקודות שבהן $y$ אינה גזירה (שימו לב ל-$x=0$).
\end{hint}

\begin{solution}
הפונקציה $y=x^{1/3}+2x^{4/3}$ רציפה בקטע הסגור $[-8,1]$, ולכן לפי \textbf{משפט ויירשטראס} היא מקבלת בו מקסימום ומינימום. נקודות אלה הן בקצות הקטע או בנקודות פנימיות קריטיות (שבהן $y'=0$ או $y'$ לא קיימת), לפי משפט פרמה.

\textbf{נגזרת} (עבור $x\ne0$):
\[ y' = \frac13x^{-2/3}+\frac83x^{1/3} = \frac{1+8x}{3\sqrt[3]{x^2}}. \]
\begin{itemize}
  \item $y'=0\iff x=-\frac18\in(-8,1)$.
  \item ב-$x=0$ הנגזרת אינה קיימת: $\frac{y(h)-y(0)}{h}=\frac{h^{1/3}(1+2h)}{h}=\frac{1+2h}{h^{2/3}}\to+\infty$. לכן $x=0$ נקודה קריטית.
\end{itemize}

\textbf{ערכי הפונקציה במועמדים:}
\begin{align*}
y(-8) &= \sqrt[3]{-8}\,(1-16) = (-2)(-15) = 30,\\
y\left(-\tfrac18\right) &= \sqrt[3]{-\tfrac18}\left(1-\tfrac14\right) = -\tfrac12\cdot\tfrac34 = -\tfrac38,\\
y(0) &= 0,\\
y(1) &= 1\cdot 3 = 3.
\end{align*}

\textbf{תשובה:} הערך הגדול ביותר הוא $30$ (מתקבל ב-$x=-8$), והערך הקטן ביותר הוא $-\frac38$ (מתקבל ב-$x=-\frac18$).
\end{solution}
